\documentclass{article}
\usepackage{amsmath}
\title{Q3P9: Contract representation in an already enforced market}
\author{Ada}
\date{}
\begin{document}
\maketitle

\section*{The pair in two sentences}
Q, \emph{Competitive Market Behavior of LLMs}, studies language-model traders in a single-unit double auction,
reporting inefficient allocations and reluctance to cross profitable spreads despite immediate execution
(ledger entries 1--3, 13). P, \emph{Commitment To Cooperation With Self-Negotiated Contracts}, studies
bargaining and navigation in Colored Trails using automatic tile transfers, interpreted natural-language
contracts, and finishing-contingent point transfers, with some contract conditions suppressing further trade
(entries 1--2, 5, 15).

\section*{The connexion pursued}
The initial proposal was to transplant P's contracts into Q to improve market performance. Both peers rejected
a generic settlement-commitment explanation: Q already executes crossing orders atomically, whereas P's
nonbinding Pay-for-Partner agreements permit later refusal, and its contract variants change different
primitives (entries 1--4). The narrower candidate asks whether, with executable obligations fixed, contractual
wording makes agents infer extra restrictions on profitable actions that remain permitted. Q supplies an
already enforced comparison; P supplies incomplete agreements requiring further coordination (entries 9, 11,
13, 24). This is a proposed experiment, not an established common mechanism.

\section*{What was found and its support}
\paragraph{Suggestive behavior, without causal isolation.}
On P's mutually dependent boards under regular trading, natural-language contracts give reported normalized
joint reward \(0.64\), versus \(0.93\) without contracts. Compared with programmatic tile contracts, they
produce similar mean executed coverage (approximately \(2.6\) moves) but fewer trades, \(1.95\) versus
\(3.28\). Q reports incremental quoting despite profitable crossing opportunities (entries 1--2, 5, 15; both
peers' notes). These observations motivate the connexion but do not identify its cause. The regular-trading
comparison must remain separate from P's delayed Pay-for-Partner results.

\paragraph{A bilateral payment can reduce to an effective price.}
Ada's reduction assumes quasilinear utility, no financing constraint or additional obligations, and a net
buyer-to-seller payment \(d\) triggered only by the same atomic trade at price \(p\). For buyer value \(v\)
and seller cost \(c\),
\[
u_B=v-p-d,\qquad u_S=p+d-c,\qquad
p_{\mathrm{eff}}=p+d,\qquad u_B+u_S=v-c.
\]
Thus the effective price reproduces the allocation and realized payoffs (entry 7; Ada's note). This is not
strategic equivalence: partner-specific payments can reorder quotes, commitments can change information or
available actions, and payments contingent on third-party trades create different payoff links (entries 7, 10,
15).

\paragraph{P's both-finish bound and the correction to unconditional optimality.}
Each player starts with fourteen own-color and two green chips. Exchanging five red for five blue gives each
nine own-color, five opposite-color, and two green chips. Any six-move shortest route uses five red/blue tiles
and the green goal, so both finish with ten chips and score \(20+5(10)=70\) each (entries 8, 10). On
asymmetric boards, exchanging five red for six blue instead gives inventories
\[
I_R=(9R,6B,2G),\qquad I_B=(5R,8B,2G).
\]
Both shortest routes remain affordable and rewards are \((75,65)\), strictly above no-interaction baselines
\((b_R,0)\) with \(b_R\leq70\) (entries 12, 15).

The precise surviving bound is conditional. Without point transfers, a nonfinisher earns zero; strict
improvement over both nonnegative baselines therefore requires both to finish. At least twelve of the initial
thirty-two chips must be consumed, giving
\[
R_{\mathrm{joint}}\leq 2(20)+5(32-12)=140.
\]
The asymmetric swap attains this bound. These are feasible constructions, not equilibrium or observed-policy
results (entries 20--22, 24).

Emmy reopened the discussion because entries 8 and 12 had incorrectly treated \(140\) as an unconditional
optimum. Under the written arbitrary-quantity exchange rule, Red can give one red chip for Blue's fourteen
blue and two green chips, leaving
\[
I_R=(13R,14B,4G),\qquad I_B=(1R,0B,0G).
\]
Red finishes with twenty-five chips, scoring \(145\); Blue fails and scores zero. If a programmatic points
contract (PPC) pays a nonfinisher, Red's prior promise of ten points upon Red's finish changes this to
\((135,10)\). Both then beat asymmetric baselines, and joint reward exceeds \(140\) (entries 17--20). This
does not Pareto-dominate \((75,65)\), since Blue receives less (entries 20, 22).

Both peers incorporated the correction. PPC may expand the set of strictly baseline-improving rewards by
compensating noncompletion, so the earlier dismissal of P's gains as necessarily behavioral was too strong.
The construction remains conditional on implementation semantics, including payments to nonfinishers, trade
quantities, scoring, and termination. P's prompt that nonfinishers receive zero total points needs
reconciliation with the PPC rule; the repository audit failed (entries 20--24; Ada's note). Emmy's finite
verifier reports \(343200\) board/path/inventory checks, supporting resource feasibility rather than
implementation fidelity (entry 21).

\paragraph{Welfare versus prices.}
Q's ranked gains are \(2.5,2,1.5,1,0.5,0\), so \(W^*=7.5\) and five positive-gain trades suffice. Emmy pairs
the five highest-value buyers with the five lowest-cost sellers at prices \((0.75,1,1.25,1.5,1.75)\): full
surplus coexists with Q's price-dispersion coefficient of approximately \(41.46\). Conversely, one trade
between value \(3.25\) and cost \(0.75\) at price \(2\) has zero dispersion but only \(1/3\) of maximum
surplus (entries 7, 12; Emmy's note). These feasible allocation/payment examples are measurement guardrails,
not equilibrium claims.

\section*{Objections and their disposition}
Equal mean executed coverage does not match contract quality, paths, inventories, timing, or residual gains;
coverage is itself downstream of treatment. P's audit of approved moves cannot bound false negatives among
rejected moves. Consequently, enforcement errors and negotiated terms remain confounded with representation
(entries 3, 5, 13). Fixed terms and deterministic enforcement are proposed controls, not completed
resolutions.

Q agents exit after one trade, so P-like residual trading by an already served agent is unavailable. The
proposed Q treatment instead re-renders unfilled standing orders while preserving their executable meaning
(entries 3, 6, 11). Refusing a currently profitable trade can also reflect rational waiting. An explicitly
disclosed final-opportunity probe removes that alternative: crossing a standing ask \(a<v\) yields \(v-a>0\),
while noncrossing yields zero. This deadline disclosure is an experimental addition (entry 11; both notes).

Finally, P's explanations that further trade is unnecessary do not establish a belief that it is forbidden.
Permission judgments and mistaken plan sufficiency require separate probes (entry 24). Attention, planning,
context burden, and interpretation remain alternative explanations; a generic prompt effect would not
establish the proposed mechanism (entries 13, 15, 24).

\section*{Sources outside Q and P}
The peers consulted Hantel's author summary at
\begin{center}
\small\texttt{agentsquared.org/research/small-bias-large-failure}
\end{center}
describing market price anchors with placebo
and debiasing controls, and Sajja et al.'s \emph{Evaluating Rational Contracting in Natural Language}
(ContractSim), sections 2 and 4, at \texttt{arxiv.org/html/2608.10475v1} (entries 5, 9, 13). Ada inspected the
author summary rather than its full linked paper; the SSRN link recorded in entry 5 was a search lead. These
sources narrow possible novelty but do not establish it.

Ada also consulted Q's repository README, at
\begin{center}
\small\texttt{github.com/jswistak/competitive-market-simulation}
\end{center}
The configuration seed controls zero-intelligence trader randomness,
while mechanism scheduling is unseeded. Paired rollouts therefore require explicit scheduling coupling or
independent market clusters (entry 10). P's anonymous repository could not be opened and its API returned
\texttt{not\_connected} (entry 20). No successful code audit is claimed.

\section*{What remains open}
The material is thin on causal evidence: neither peer ran an LLM experiment or market replication. The record
establishes economic boundaries, a corrected conditional reward construction, and a research design. Pinned
code semantics, equivalent displays, separation of scope mistakes from attention or planning, full-market
effect sizes and power, and broader novelty remain unresolved (entries 13, 15, 21, 24). Joint reward,
individual baseline improvement, both-finish rates, and price dispersion must be reported separately.

\section*{The smallest next step}
Audit a terminal Q snapshot in a pinned implementation, then run the paired display diagnostic from entry 11:
identical order state and permissions, rendered as an order record or conditional obligation. On separate
cloned contexts, measure permission errors and compare an entailed permission clarification with an
equal-budget generic reminder. A display-dependent false prohibition, costly failure to cross, and selective
rescue would support proceeding to full-market tests; changed action rates alone would support only framing
sensitivity (entries 11, 13, 24; both notes). This settles the first diagnostic question, not market
significance or novelty. Separately, the smallest check of P's corrected bound is to replay the explicit
exchange and ten-point PPC in its implementation and inspect the nonfinisher's final reward (entries 20--22).

\end{document}
